% !TEX program = xelatex
% ===========================================================
% Yanda Cheng -- 西门子CT临床科研支持专员 双语简历
% Siemens Healthineers CT Clinical Research Support -- Bilingual Resume
%
% 编译方式 / Compile with: XeLaTeX (必须 / required)
% 中文版 Chinese version: 第 1-2 页 / Pages 1-2
% 英文版 English version: 第 3-4 页 / Pages 3-4
% 中文字体 Chinese font: Microsoft YaHei (Windows 内置 / built-in)
% Overleaf 用户请将 Microsoft YaHei 替换为 Noto Sans CJK SC
% ===========================================================

\documentclass[letterpaper,10pt]{article}

% ---------- Packages ----------
\usepackage{iftex}
\usepackage[empty]{fullpage}
\usepackage{titlesec}
\usepackage[usenames,dvipsnames]{color}
\usepackage{enumitem}
\setlist[itemize]{noitemsep, topsep=3pt, left=1.5em}
\usepackage[hidelinks]{hyperref}
\usepackage{fancyhdr}
\usepackage{tabularx}

% ---------- Engine compatibility ----------
% This resume contains Chinese and requires either XeLaTeX or LuaLaTeX.
% pdfLaTeX will NOT work.
\ifPDFTeX
  \errmessage{This file requires XeLaTeX or LuaLaTeX (not pdfLaTeX). Please change the compiler to XeLaTeX or LuaLaTeX.}
\fi

\ifXeTeX
  \usepackage{fontspec}
  \usepackage{xeCJK}
\fi
\ifLuaTeX
  \usepackage{fontspec}
  \usepackage{luatexja}
  \usepackage{luatexja-fontspec}
\fi

% ---------- Fonts ----------
% Prefer Windows font locally; fall back to Noto on Overleaf.
\ifXeTeX
  \IfFontExistsTF{Microsoft YaHei}{
    \setCJKmainfont{Microsoft YaHei}
    \setCJKsansfont{Microsoft YaHei}
  }{
    \setCJKmainfont{Noto Sans CJK SC}
    \setCJKsansfont{Noto Sans CJK SC}
  }
\fi
\ifLuaTeX
  \IfFontExistsTF{Microsoft YaHei}{
    \setmainjfont{Microsoft YaHei}
  }{
    \setmainjfont{Noto Sans CJK SC}
  }
\fi
\setmainfont{Arial}
\linespread{1.03}\selectfont

% ---------- Page Style ----------
\pagestyle{fancy}
\fancyhf{}
\renewcommand{\headrulewidth}{0pt}
\renewcommand{\footrulewidth}{0pt}

% ---------- Margins ----------
\addtolength{\oddsidemargin}{-0.5in}
\addtolength{\textwidth}{1in}
\addtolength{\topmargin}{-0.9in}
\addtolength{\textheight}{1.15in}

% ---------- Section Formatting ----------
% Note: \scshape removed for CJK compatibility; bold underline retained
\titleformat{\section}{\large\bfseries}{}{0em}{}[\titlerule]
\titlespacing{\section}{0pt}{5pt plus 1pt minus 1pt}{3pt}

% ---------- Custom Commands ----------
\newcommand{\resumeItem}[1]{\item\small{#1}}

\newcommand{\resumeSubheading}[4]{
  \vspace{-1pt}\item
  \begin{tabular*}{0.97\textwidth}[t]{l@{\extracolsep{\fill}}r}
    \textbf{#1} & \small#2 \\
    \textit{\small#3} & \textit{\small#4} \\
  \end{tabular*}\vspace{-2pt}
}

\newcommand{\projheading}[2]{
  \vspace{1pt}\item
  \begin{tabular*}{0.97\textwidth}[t]{l@{\extracolsep{\fill}}r}
    \textbf{#1} & \textit{\small#2} \\
  \end{tabular*}\vspace{2pt}
}

\newcommand{\resumeSubHeadingList}{\begin{itemize}[leftmargin=0.0in, label={}]}
\newcommand{\resumeSubHeadingListEnd}{\end{itemize}}
\newcommand{\resumeItemListStart}{\begin{itemize}[leftmargin=0.2in]}
\newcommand{\resumeItemListEnd}{\end{itemize}}
\newcommand{\projItemListStart}{\begin{itemize}[leftmargin=0.2in, topsep=3pt]}
\newcommand{\projItemListEnd}{\end{itemize}\vspace{2pt}}

% =================================================================
\begin{document}
% =================================================================


% #################################################################
% #################### 中文版 CHINESE VERSION #####################
% #################################################################

% ---------- 头部 Header ----------
\begin{center}
  {\huge \textbf{程彦达 \quad Yanda Cheng}} \\
  \vspace{4pt}
  \small
  +1-859-338-3379 \quad$|$\quad
  \href{mailto:nickcp39@gmail.com}{nickcp39@gmail.com} \quad$|$\quad
  \href{https://yanda-cheng.com}{yanda-cheng.com} \quad$|$\quad
  \href{https://www.linkedin.com/in/yanda-cheng-4888a216b}{linkedin.com/in/yanda-cheng-4888a216b} \quad$|$\quad
  \href{https://scholar.google.com/citations?view_op=search_authors&mauthors=Yanda+Cheng+University+at+Buffalo}{Google Scholar Profile}
\end{center}

\vspace{-2pt}

% ---------- 教育背景 ----------
\section{教育背景}

\begin{itemize}[leftmargin=0.0in, label={}, itemsep=3pt, topsep=2pt]
  \item
    \textbf{生物医学工程 博士（在读）}\hfill
    \textit{布法罗大学（State University of New York at Buffalo）}\hfill
    GPA: 4.0/4.0 \hfill 2021--2026（预计5月）
  \item
    \textbf{生物医学工程 硕士}\hfill
    \textit{康奈尔大学（Cornell University，常春藤盟校）}\hfill
    2020--2021
  \item
    \textbf{电子与计算机工程 学士}\hfill
    \textit{肯塔基大学（University of Kentucky）}\hfill
    2015--2020
\end{itemize}

\vspace{-2pt}

% ---------- 工作经历 ----------
\section{工作经历}

\begin{itemize}[leftmargin=0in, label={}]

  % --- Seno Medical ---
  \item
  \textbf{机器学习工程师实习 · Seno Medical（医疗成像设备公司）}\hfill 纽约 \\
  \textit{临床数据集构建 · AI辅助医学影像评估 · 合规部署} \hfill \textit{2024年5月 -- 8月}
  \begin{itemize}[leftmargin=0.2in, topsep=1pt]
    \item 与临床团队协作，构建 2000+ 条医学影像 QA 数据集并制定标注规范，完成 LLaMA（LoRA）微调与检索增强（RAG）系统开发，用于医生端问答与报告辅助。
    \item 设计端到端评估框架（标签体系、自动化指标与误差分析脚本），将 hallucination 率从约 30\% 降低至 15\%，并将 Recall@5 提升至约 90\%，形成可复现的评估流程。
    \item 搭建模块化 RAG 后端：完成文本切分与摄取任务、基于 FAISS 的向量库、检索与重排序层以及推理服务编排，支持后续模型快速迭代与对比实验。
    \item 设计符合 PHI 法规的本地化安全部署方案（FastAPI\,+\,Docker），并为临床团队输出部署文档与模型评估报告，便于在医院内部网络长期运行与维护。
  \end{itemize}

  \vspace{3pt}

  % --- Roswell Park ---
  \item
  \textbf{研究科学家 · Roswell Park 综合癌症中心}\hfill 纽约州布法罗 \\
  \textit{医学影像算法研发 · 临床数据分析} \hfill \textit{2023年5月 -- 8月}
  \begin{itemize}[leftmargin=0.2in, topsep=1pt]
    \item 参与肿瘤相关多模态医学影像研究，负责图像预处理、特征工程与模型评估，保障数据质量与实验可重复性。
    \item 与放射科、肿瘤科医生协作开展数据整理与结果验证，将复杂算法结果转化为临床可解释的影像与统计指标，支撑论文撰写与会议展示。
    \item 搭建影像+结构化报告的数据流水线（去标识化、模式规范化、自动质控），为后续生成式建模与预测模型提供标准化数据资产。
    \item 在 AWS 环境中完成模型训练与调参，建立实验追踪脚本与对比基线，缩短不同模型方案之间的迭代周期。
  \end{itemize}

  \vspace{3pt}

  % --- HydroSense ---
  \item
  \textbf{联合创始人 · HydroSense Tech LLC（高精度传感器公司）}\hfill 北京 / 新加坡 \\
  \textit{传感器硬件研发 · 嵌入式ML · IoT数据平台} \hfill \textit{2015年 -- 至今（兼职）}
  \begin{itemize}[leftmargin=0.2in, topsep=1pt]
    \item 联合创立物联网高精度传感创业公司，负责核心算法与系统架构设计，推动产品从试点走向规模化应用，年营收达到百万美元级别，落地于 5+ 国家市政与工业场景。
    \item 设计传感器信号链（前端放大、采样同步、数字滤波），在显著温度与电磁干扰条件下实现稳定测量，为后续 AI 校准与远程监控打下基础。
    \item 构建基于 Python / MLflow 的 AI 驱动校准与预测框架，以回归与神经网络模型建模非线性与温度依赖，将现场漂移降低约 30\%，并实现设备健康状态长期追踪。
    \item 搭建支持 OTA 的远程监控与自诊断系统（C++ 固件 + 网关服务），实现 5000+ 设备的集中配置管理与异常预警，减少现场维护频次与成本。
    \item 持有多项传感器算法与硬件发明专利（信号处理、漂移补偿、EMI 抑制等，2017--2023），将研究成果沉淀为可保护的知识产权。
  \end{itemize}

\end{itemize}

\vspace{-2pt}

% ---------- 代表性科研项目 ----------
\section{代表性科研项目（与CT影像临床科研支持高度相关）}

\begin{itemize}[leftmargin=0in, label={}]

  \projheading{慢性足部溃疡纵向光声影像与预后预测建模}{npj Imaging, 2026}
  \projItemListStart
    \item 设计纵向随访方案，纳入 N\,>\,400 例糖尿病足溃疡患者；协调内分泌科、血管外科、伤口护理科完成 IRB 申请与多中心数据采集
    \item 融合影像组学特征与 LASSO / 随机森林预测模型，AUC\,>\,0.85，实现溃疡愈合概率定量预测，支撑临床决策
    \item \textit{可迁移至：}CT灌注成像纵向随访研究、低剂量CT肺癌筛查队列预测建模
  \projItemListEnd

  \projheading{乳腺光声-超声多模态自动化成像系统（OneTouch）}{IEEE Transactions on Medical Imaging, 2025（被引8次）}
  \projItemListStart
    \item 参与多模态临床成像系统研发，结合 AI 辅助分析，面向乳腺癌早期筛查；与放射科及外科医生合作设计成像协议、完成临床验证
    \item \textit{可迁移至：}西门子CT联合超声多模态乳腺筛查联合研究
  \projItemListEnd

  \projheading{光声图像无监督 AI 去噪（Noise2Noise Network）}{Biomedical Optics Express, 2024（被引14次）}
  \projItemListStart
    \item 提出无需配对干净图像的深度学习去噪框架，显著提升低信噪比医学图像质量
    \item \textit{可迁移至：}\textbf{低剂量CT图像AI去噪}（SOMATOM低剂量肺癌筛查研究方向，图像质量提升是热门发表点）
  \projItemListEnd

  \projheading{足部血管双扫描光声断层成像（Dual-Scan PA Tomography）}{IEEE Transactions on UFFC, 2023（被引21次）}
  \projItemListStart
    \item 开发三维定量血管成像协议，用于外周动脉疾病（PAD）评估；与血管外科团队制定患者招募与成像标准
    \item \textit{可迁移至：}西门子CT血管造影（CTA）研究、下肢PAD定量分析
  \projItemListEnd

  \projheading{吞咽困难光声影像评估（婴儿猪活体/离体模型）}{Med-X, 2025}
  \projItemListStart
    \item 首次将光声影像技术应用于吞咽功能评估；与言语治疗科合作，主导实验设计、影像处理与论文撰写全过程
    \item \textit{展示：}开拓全新临床应用场景并推动发表的科研组织能力
  \projItemListEnd

  \projheading{乳腺癌早期检测转化研究与创业}{布法罗大学创业大赛获奖，2023--至今}
  \projItemListStart
    \item 主导乳腺癌光声早期检测的转化项目，与布法罗本地医院合作设计患者研究，完成创业路演并获批种子基金 \$3,000
    \item \textit{展示：}科研成果转化与产业对接能力
  \projItemListEnd

\end{itemize}

\vspace{-2pt}

% ---------- 代表性论文 ----------
\section{代表性论文（第一 / 共同作者，同行评审，已发表）}

\begin{itemize}[leftmargin=0.2in, itemsep=1pt]
  \item \small \textbf{Cheng Y.}, Huang C., et al. Predictive modeling of chronic foot ulcer outcomes using longitudinal photoacoustic imaging. \textit{npj Imaging} 4(1), 12, \textbf{2026}.
  \item \small \textbf{Cheng Y.}, Huang C., Bing R.W., et al. Dysphagia assessment based on photoacoustic imaging: a pilot ex vivo and in vivo study in infant swine. \textit{Med-X} 3(1), 1--9, \textbf{2025}.
  \item \small \textbf{Cheng Y.}, Zheng W., Bing R.W., et al. Unsupervised denoising of photoacoustic images based on the Noise2Noise network. \textit{Biomedical Optics Express} 15(8), 4390--4405, \textbf{2024}.\quad\textbf{[被引 14 次]}
  \item \small \textbf{Cheng Y.}, Huang C., et al. Dual-scan photoacoustic tomography for the imaging of vascular structure on foot. \textit{IEEE Trans. Ultrason. Ferroelectr. Freq. Control} 70, \textbf{2023}.\quad\textbf{[被引 21 次]}
  \item \small Huang C., \textbf{Cheng Y.}, et al. Enhanced clinical photoacoustic vascular imaging through a skin localization network and adaptive weighting. \textit{Photoacoustics} 42, 100690, \textbf{2025}. [被引 4 次]
  \item \small Zhang H., Zheng E., ..., \textbf{Cheng Y.}, et al. OneTouch automated photoacoustic and ultrasound imaging of breast in standing pose. \textit{IEEE Trans. Medical Imaging}, \textbf{2025}. [被引 8 次]
  \item \small Liu X., \textbf{Cheng Y.}, et al. Simultaneous measurements of tissue blood flow and oxygenation using a wearable fiber-free optical sensor. \textit{J. Biomed. Optics} 26(1), 012705, \textbf{2021}.\quad\textbf{[被引 38 次]}
\end{itemize}
\vspace{2pt}
\noindent\small\textbf{总引用次数：约100次（Google Scholar）\quad 共11篇 SCI 论文\quad 期刊：npj Imaging, IEEE TMI, IEEE TUFFC, Photoacoustics, BOE, Med-X, JBO}

\vspace{-2pt}

% ---------- 专业技能 ----------
\section{专业技能}
\noindent
\small
\textbf{医学影像处理：} 图像重建、去噪（深度学习/传统）、配准、分割、3D可视化；CT、光声、超声、多光子显微成像 \\
\textbf{AI / 机器学习：} CNN, UNet, Transformer, LSTM, 影像组学（Radiomics），LASSO, SVM, 随机森林，深度学习去噪 \\
\textbf{临床数据分析：} 纵向分析、生存分析（Kaplan-Meier）、ROC/AUC, 线性/逻辑回归，统计检验（Python / R / SPSS）\\
\textbf{科研协作流程：} IRB/伦理委员会申请、多中心患者招募、数据质量控制、SCI论文撰写与投稿、学术报告\\
\textbf{编程工具：} Python, C/C++, MATLAB, R, SQL；PyTorch, scikit-learn, OpenCV, Pandas, Matplotlib, Docker \\
\textbf{语言能力：} 普通话（母语），英语（流利——学术写作、科研汇报、跨国团队协作）

\vspace{-2pt}

% ---------- 荣誉获奖 ----------
\section{荣誉与获奖}
\begin{itemize}[leftmargin=0.2in, itemsep=1pt]
  \item \small 布法罗大学创业竞赛获奖，获批创业种子基金 \textbf{\$3,000}（乳腺癌光声早期检测创业项目，2023）
  \item \small NIH 资助科研项目主要参与者（PI: Prof. Jun Xia，UB，持续资助中）
  \item \small 布法罗大学优秀助教（BME 503 图像处理 / BME 302 医疗器械）
\end{itemize}


% #################################################################
% ################### 英文版 ENGLISH VERSION ######################
% #################################################################

\clearpage

% ---------- Header ----------
\begin{center}
  {\huge \textbf{Yanda Cheng}} \\
  \vspace{4pt}
  \small
  +1-859-338-3379 \quad$|$\quad
  \href{mailto:nickcp39@gmail.com}{nickcp39@gmail.com} \quad$|$\quad
  \href{https://yanda-cheng.com}{yanda-cheng.com} \quad$|$\quad
  \href{https://www.linkedin.com/in/yanda-cheng-4888a216b}{linkedin.com/in/yanda-cheng-4888a216b} \quad$|$\quad
  \href{https://scholar.google.com/citations?view_op=search_authors&mauthors=Yanda+Cheng+University+at+Buffalo}{Google Scholar Profile}
\end{center}

\begin{center}
  \small
  \textbf{Target Role:} CT Clinical Research Support Specialist $\cdot$ Siemens Healthineers $\cdot$ Shanghai $\cdot$
  Available May 2026 (PhD graduation)
\end{center}

\vspace{-2pt}

% ---------- Education ----------
\section{Education}

\begin{itemize}[leftmargin=0.0in, label={}, itemsep=2pt, topsep=2pt]
  \item \textbf{PhD} in Biomedical Engineering \hfill \textit{University at Buffalo (SUNY), Buffalo, NY} \hfill GPA: 4.0/4.0 \hfill 2021--2026 (Expected)
  \item \textbf{MEng} in Biomedical Engineering \hfill \textit{Cornell University (Ivy League), Ithaca, NY} \hfill 2020--2021
  \item \textbf{BS} in Electrical \& Computer Engineering \hfill \textit{University of Kentucky, Lexington, KY} \hfill 2015--2020
\end{itemize}

\vspace{-2pt}

% ---------- Work Experience ----------
\section{Work Experience}

\begin{itemize}[leftmargin=0in, label={}]

  % --- Seno Medical ---
  \item
  \textbf{Applied ML Engineer Intern -- Seno Medical (Medical Imaging Device Company)}\hfill New York \\
  \textit{Clinical Dataset Development, AI-Assisted Medical Imaging, PHI-Compliant Deployment} \hfill \textit{May 2024 -- Aug 2024}
  \begin{itemize}[leftmargin=0.2in, topsep=1pt]
    \item Collaborated with clinicians to build a 2k+ medical imaging QA dataset with clear annotation guidelines, and fine-tuned a LLaMA-based model (LoRA) as the core of a retrieval-augmented QA system.
    \item Designed an end-to-end evaluation harness with label schema (factuality, coverage, style), automatic metrics, and Python error-analysis scripts, reducing hallucination from \textasciitilde30\% to 15\% and increasing Recall@5 to 90\%.
    \item Implemented a modular RAG backend with text chunking/ingestion jobs, a FAISS-based vector store, retriever/re-ranker layer, and LLM orchestration, enabling reproducible experiments across model and retrieval variants.
    \item Engineered PHI-compliant on-prem AI deployment with FastAPI and Docker for hospital-network environments, and produced clinician-facing reports summarizing model behavior, limitations, and integration recommendations.
  \end{itemize}

  \vspace{3pt}

  % --- Roswell Park ---
  \item
  \textbf{Research Scientist -- Roswell Park Comprehensive Cancer Center}\hfill Buffalo, NY \\
  \textit{Medical Imaging Algorithm Development, Clinical Data Analysis} \hfill \textit{May 2023 -- Aug 2023}
  \begin{itemize}[leftmargin=0.2in, topsep=1pt]
    \item Designed a multimodal data pipeline aligning imaging with structured text reports, including de-identification, schema normalization, and automated QC, to produce clean datasets for downstream modeling.
    \item Built dataset generation jobs (filtering, sampling, augmentation, metadata tagging) to construct balanced image--report datasets for conditional generation and predictive modeling.
    \item Implemented a domain-specific latent diffusion model (Stable Diffusion--style, DDIM/DPMS sampling) on AWS for conditional image-to-text report generation and synthetic data augmentation.
    \item Automated GPU training and experiment tracking with standardized configs, logging, and comparison scripts, enabling fast iteration and reproducible baselines across model variants.
  \end{itemize}

  \vspace{3pt}

  % --- HydroSense ---
  \item
  \textbf{Co-founder \& Core Engineer -- HydroSense Tech LLC}\hfill Beijing / Singapore \\
  \textit{IoT Sensing, Embedded ML, Signal Processing} \hfill \textit{2015 -- Present (part-time)}
  \begin{itemize}[leftmargin=0.2in, topsep=1pt]
    \item Co-founded a precision sensing startup and led full-stack development (optics, embedded firmware, data intelligence), scaling deployments from pilots to multi-million-dollar annual revenue across municipal and industrial customers in 5+ countries.
    \item Designed the end-to-end sensor signal pipeline, including photodiode amplification, sampling synchronization, and fixed-point digital filtering, achieving sub-percent repeatability under temperature and EMI drift.
    \item Built an AI-driven calibration and prediction framework (Python/MLflow) using regression and neural compensation models to capture nonlinearities and temperature dependence, reducing field drift by \textasciitilde30\% while tracking long-term device health.
    \item Developed a self-diagnostic telemetry and remote-monitoring system (C++ firmware + Python services) with OTA updates and anomaly alerts, enabling proactive maintenance and significantly reducing on-site service visits for a 5000+ device fleet.
    \item \textbf{Patents:} Co-inventor on multiple granted patents covering rainfall sensing hardware, adaptive signal processing, and intelligent embedded calibration algorithms (2017--2023).
  \end{itemize}

\end{itemize}

\vspace{-2pt}

% ---------- Research Projects ----------
\section{Selected Research Projects (Directly Relevant to CT Imaging Research Support)}

\begin{itemize}[leftmargin=0in, label={}]

  \projheading{Longitudinal PA Imaging \& Predictive Modeling of Chronic Foot Ulcers}{npj Imaging, 2026}
  \projItemListStart
    \item Designed full longitudinal follow-up protocol for N\,$>$\,400 diabetic foot ulcer patients; coordinated \textbf{IRB approval}, multi-center patient enrollment, and data collection with endocrinology, vascular surgery, and wound care teams
    \item Combined radiomics features with LASSO/Random Forest models achieving AUC\,$>$\,0.85 for quantitative ulcer healing prediction, directly supporting clinical decision-making
    \item \textit{Transferable to:} CT perfusion longitudinal studies, low-dose CT lung cancer screening cohort risk stratification
  \projItemListEnd

  \projheading{OneTouch Multi-modal Breast PA/US Automated Imaging System}{IEEE Trans. Medical Imaging, 2025 (cited 8$\times$)}
  \projItemListStart
    \item Co-developed a multi-modal clinical imaging platform integrating AI-assisted diagnosis for early breast cancer screening; coordinated with radiologists and surgeons on clinical protocol design and image interpretation validation
    \item \textit{Transferable to:} Siemens CT + US multi-modal breast cancer screening collaborative research
  \projItemListEnd

  \projheading{Unsupervised AI Denoising for Medical Images (Noise2Noise Network)}{Biomedical Optics Express, 2024 (cited 14$\times$)}
  \projItemListStart
    \item Proposed deep learning denoising framework requiring no paired clean images; validated on photoacoustic imaging with significant SNR improvement
    \item \textit{Directly transferable to:} \textbf{Low-dose CT image AI denoising} -- a key publication area for Siemens SOMATOM low-dose lung cancer screening programs
  \projItemListEnd

  \projheading{Dual-Scan Photoacoustic Tomography for Foot Vascular Imaging}{IEEE Trans. UFFC, 2023 (cited 21$\times$)}
  \projItemListStart
    \item Developed quantitative 3D vascular imaging protocol for peripheral arterial disease (PAD) assessment; established clinical imaging standards in collaboration with a vascular surgery team
    \item \textit{Transferable to:} Siemens CT angiography (CTA) for lower extremity PAD quantitative analysis
  \projItemListEnd

  \projheading{Dysphagia Assessment via Photoacoustic Imaging (Infant Swine Model)}{Med-X, 2025}
  \projItemListStart
    \item First application of photoacoustic imaging for swallowing function evaluation; collaborated with speech-language pathology clinicians on experimental design, image processing, and full manuscript preparation
    \item \textit{Demonstrates:} Ability to open new clinical application domains and deliver complete research output
  \projItemListEnd

  \projheading{Translational Research \& Startup for PA-Based Breast Cancer Detection}{UB Entrepreneurship Award, 2023--Present}
  \projItemListStart
    \item Led translational research project for early breast cancer detection; partnered with local hospitals to design patient studies and validate clinical feasibility; secured \textbf{\$3,000 startup grant} and entrepreneurship prizes
    \item \textit{Demonstrates:} Clinical translation thinking and stakeholder relationship building (directly applicable to KOL management in a medical device company)
  \projItemListEnd

\end{itemize}

\vspace{-2pt}

% ---------- Publications ----------
\section{Selected Publications (First / Co-Author, Peer-Reviewed)}

\begin{itemize}[leftmargin=0.2in, itemsep=1pt]
  \item \small \textbf{Cheng Y.}, Huang C., et al. Predictive modeling of chronic foot ulcer outcomes using longitudinal photoacoustic imaging. \textit{npj Imaging} 4(1), 12, \textbf{2026}.
  \item \small \textbf{Cheng Y.}, Huang C., Bing R.W., et al. Dysphagia assessment based on photoacoustic imaging: a pilot ex vivo and in vivo study in infant swine. \textit{Med-X} 3(1), 1--9, \textbf{2025}.
  \item \small \textbf{Cheng Y.}, Zheng W., et al. Unsupervised denoising of photoacoustic images based on the Noise2Noise network. \textit{Biomed.\ Optics Express} 15(8), 4390--4405, \textbf{2024}.\quad\textbf{[14 citations]}
  \item \small \textbf{Cheng Y.}, Huang C., et al. Dual-scan photoacoustic tomography for the imaging of vascular structure on foot. \textit{IEEE Trans.\ UFFC} 70, \textbf{2023}.\quad\textbf{[21 citations]}
  \item \small Huang C., \textbf{Cheng Y.}, et al. Enhanced clinical photoacoustic vascular imaging through skin localization and adaptive weighting. \textit{Photoacoustics} 42, 100690, \textbf{2025}. [4 citations]
  \item \small Zhang H., ..., \textbf{Cheng Y.}, et al. OneTouch automated PA/US breast imaging in standing pose. \textit{IEEE Trans.\ Medical Imaging}, \textbf{2025}. [8 citations]
  \item \small Liu X., \textbf{Cheng Y.}, et al. Simultaneous tissue blood flow and oxygenation with a wearable fiber-free optical sensor. \textit{J.\ Biomed.\ Optics} 26(1), 012705, \textbf{2021}.\quad\textbf{[38 citations]}
\end{itemize}
\vspace{2pt}
\noindent\small\textbf{Total: $\sim$100 citations (Google Scholar) $|$ 11 SCI papers $|$ Venues: npj Imaging, IEEE TMI, IEEE TUFFC, Photoacoustics, BOE, Med-X, JBO}

\vspace{-2pt}

% ---------- Technical Skills ----------
\section{Technical Skills}
\noindent
\small
\textbf{Medical Imaging:} Image reconstruction, denoising (DL/classical), registration, segmentation, 3D visualization; CT, PA, US, multiphoton microscopy \\
\textbf{AI / Machine Learning:} CNN, UNet, Transformer, LSTM, Radiomics, LASSO, SVM, Random Forest, deep learning denoising \\
\textbf{Clinical Data Analysis:} Longitudinal analysis, Kaplan-Meier survival analysis, ROC/AUC, regression, statistical testing (Python\,/\,R\,/\,SPSS) \\
\textbf{Research Workflow:} IRB/ethics protocol submission, multi-site data management, SCI manuscript preparation \& submission, academic presentation \\
\textbf{Programming:} Python, C/C++, MATLAB, R, SQL; PyTorch, scikit-learn, OpenCV, Pandas, Docker \\
\textbf{Languages:} Mandarin Chinese (native), English (fluent -- academic writing, cross-national team collaboration)

\vspace{-2pt}

% ---------- Awards ----------
\section{Awards \& Honors}
\begin{itemize}[leftmargin=0.2in, itemsep=1pt]
  \item \small \textbf{\$3,000 Startup Seed Grant}, University at Buffalo Entrepreneurship Award -- PA-based Breast Cancer Detection (2023)
  \item \small Key Contributor, NIH-Funded Research Project (PI: Prof.\ Jun Xia, University at Buffalo; ongoing)
  \item \small Outstanding Teaching Assistant -- BME 503 Image Processing \& BME 302 Medical Devices, University at Buffalo
\end{itemize}


\end{document}
